% Generated by code/interviews-1/03-speed-drills/python/gen_drill.py. Do not edit. French copy by draft/pilot-quant-fr/ivgen_fr.py.
\ivfamily{Additionner et soustraire de grands nombres}{2, 6, 11, 15, 17, 28, 30, 47}{Arrondir un nombre puis corriger~: $8142 - 3678 = 8142 - 3700 + 22 = 4464$.}
\ivfamily{Deux chiffres fois deux chiffres}{13, 16, 18, 19, 24, 29, 37, 39, 45, 50}{Décomposer un facteur, comme dans $37 \times 46 = 37 \times 40 + 37 \times 6 = 1480 + 222 = 1702$. Si les facteurs sont proches, utiliser plutôt la différence de deux carrés.}
\ivfamily{Carrés}{12, 20, 21, 36, 40, 48}{$n^2 = (n - d)(n + d) + d^2$, avec $d$ choisi pour rendre un facteur rond~: $48^2 = 50 \times 46 + 4 = 2304$. Un nombre finissant par 5~: $65^2$ vaut $6 \times 7$ suivi de 25.}
\ivfamily{Autour d'un nombre rond}{9, 31, 32, 35, 42}{$103 \times 97 = 100^2 - 3^2$, et $998 \times 6 = 6000 - 12$.}
\ivfamily{Pourcentages de n'importe quel nombre}{3, 4, 7, 25, 27, 34, 43}{Prendre 1~\% ou 10~\% et monter, ou permuter les deux~: $24\%$ de 350 égale $350\%$ de 24, soit 84.}
\ivfamily{Fractions}{1, 14, 22, 23, 33, 38, 49}{Utiliser le plus petit dénominateur commun. Pour diviser, multiplier par la fraction inverse, et simplifier d'abord.}
\ivfamily{Réponses décimales}{5, 8, 10, 26, 41, 44, 46}{Connaître les huitièmes et les seizièmes~: $\tfrac18 = 0.125$ et $\tfrac1{16} = 0.0625$. Pour diviser par 40, diviser par 4 puis par 10.}
